%\needspace{4\baselineskip}
\color{uniblue}{\subsection{Орган выявления БНТ ВВ}}
\color{black}

\begin{enumerate}[label=\arabic{section}.\arabic{subsection}.\arabic{enumi}, labelsep=4pt, leftmargin=0pt, itemindent=57pt, itemsep=0pt, parsep=0pt] % если что можно крутить itemsep


\item
Для обеспечения несрабатывания «МТЗ ВВ», «ТЗНП ВВ», «ТЗНП НВ», «ТЗОП ВВ» при броске тока 
намагничивания используется блокировка по содержанию 2 гармоники в токе ВН. Блокировка основана 
на контроле отношения действующего значения второй гармонике к первой тока фаз. Имеется 
возможность регулировать порог срабатывания блокировки, с помощью уставки \textbf{«Ih2/Ih1 БНТ»}. 
Блокировка по 2 гармонике выводится из действия при величине тока ниже 0,04Iном. Блокировка по 2 
гармонике выводится из действия при величине тока ниже 0,04Iном, при срабатывании «Iср БНТ», а так 
же при выводе программной накладкой «БНТ ВВ».

\begin{figure}[H]
\centering
\includegraphics[width=1\textwidth,height=1\textheight,keepaspectratio]{\fbpath/БНТ/img/BNT.pdf}
\captionof{figure}{Функциональный блок БНТ}\label{pict:BNT}
\end{figure}




\small
\setlength{\extrarowheight}{0.05cm} % высота ячеек в таблице
\begin{longtable}
{|>{\centering\arraybackslash}m{4.5cm}|>{\centering\arraybackslash}m{4.5cm}|>{\centering\arraybackslash}m{3cm}|>{\centering\arraybackslash}m{3.5cm}|}
\caption{Перечень уставочных параметров функции БНТ\hfill\vspace{-0.5\baselineskip}}\label{ksn:tbl2} \\ 
\hhline{|-|-|-|-|}
\rowcolor{gray!30}
Наименование уставки & Диапазон & Шаг &  Значение по умолчанию \\ 
\hhline{|-|-|-|-|}
\endfirsthead

\caption*{\hspace{3pt}\emph{Продолжение таблицы \ref{ksn:tbl2}\hfill\vspace{-0.5\baselineskip}}} \\ % сделано по ГОСТ 2.105 п.6.8.7
\hhline{|-|-|-|-|}
\rowcolor{gray!30}
Наименование уставки & Диапазон & Шаг &  Значение по умолчанию \\  
\hhline{|-|-|-|-|}
\endhead
\hhline{|-|-|-|-|}

\endfoot
\endlastfoot

%%%%%%%%%%%%содержимое%%%%%%%%%%%%%%

\end{longtable}
\setlength{\extrarowheight}{0.0cm} % возврат параметра высота ячеек в таблице
\normalsize












































\end{enumerate} 